\scanpage{21}\mspage{I}
Le théorème de Jordan et l'existence de trois bons sommets pour un polygone \tikz[baseline=(a.base)]{\node[inner sep=1pt] (a) {de $\mathbb{R}^{2}$.};\draw[thin,gray,rotate around={8:(a.center)}] (a.center) ellipse (0.65cm and 0.30cm);}\inserted{plan}

\hspace*{\fill}Carlos C.~\unsure{Contou|Coutou}-Carrère\hspace*{\fill}

\margin{$\bullet$}On se propose de prouver les deux théorèmes suivantes :

\uline{Théorème} 1 (Jordan). Tout polygone $\Gamma$ est un $J$-polygone (cf. def 1)\margin{\tikz{\draw[gray,thin] (.1,.6) .. controls (-.2,.9) and (-.15,.1) .. (.25,-.3);\draw[gray,thin] (.25,-.35) .. controls (.35,-.6) and (.5,-.8) .. (.7,-.9);}}

\uline{Théorème} 2 (Dehn (?) non-publié). Si $\Gamma$ est un polygone, il existe au moins trois bons sommets (cf. def.\hspace{1em}) dont deux au moins sont tels que les triangles correspondants soient \unsure{interieurs|inteieurs|intérieurs}, i.e. contenus dans $\mathrm{int}_{\Gamma}\cup\Gamma$ (cf. def 2).

1- \uline{Preliminaires}

\margin{$\bullet\ \bullet$}On entend par polygone la réalisation \uline{géom}étrique dans \uline{$\mathbb{R}^{2}$} \margin{\tikz{\draw[gray,thin] (0,0) .. controls (.6,.35) and (1.15,.38) .. (1.5,.2);\draw[gray,thin] (.1,-.05) .. controls (.8,-.15) and (1.55,-.3) .. (2,-.7);}}d'un \margin{$\bullet$}polygo\uline{ne abs}trait, connexe\inserted{et non vide}. Si $\Gamma$ est un polygone on dénote par \margin{$\bullet$}$S(\Gamma)$ (resp. $A(\Gamma)$) l'ensemble des sommets (resp. arêtes) gé\uline{ometriq}\unsure{ues} de $\Gamma$.

\margin{$\cdot\quad\cdot$}Le terme ar\uline{c de} polygone désigne la réalisation \uline{géo}metrique dans $\mathbb{R}^{2}$ d'un 1-complexe connexe$\langle$ \margin{or. ?}tel que chaque sommet soit incident à 1 ou à 2 arêtes, et tel qu'il existe exactement 2 sommets qui sont incidents à une seule arête (pas néc. la même).

On appelle extrémités \unsure{d'eun|d'un|deun} arc de polygone les points de $\mathbb{R}^{2}$ qui correspondent à ces deux sommets. Si $\Gamma_{xy}$ est un arc de polygone d'\unsure{extremités|extrémités} $x,y$ on dénote par $S(\Gamma_{xy})$ (resp. $A(\Gamma_{xy})$) l'ensemble de ses sommets (resp. arêtes) géometriques.

\margin{$\bullet$}Par $\Gamma$ (resp. $\Gamma_{xy}$) on dénote a\uline{us}si l'ensemble \inserted{de $\mathbb{R}^{2}$} sous jacent à $\Gamma$ (resp. $\Gamma_{xy}$) si $\Gamma$ \deleted{est} (resp. $\Gamma_{xy}$) est un polygone (resp. arc de polygone).

\tikz[baseline=(r.base)]{\node[inner sep=1pt] (r) {\uline{Remarque 1}};\draw[dashed,gray,line width=.2pt] ([xshift=-1pt,yshift=1pt]r.north east)..controls ([yshift=5pt]r.north) and ([xshift=-3pt,yshift=4pt]r.north west)..([xshift=-3pt]r.west)..controls ([xshift=-3pt,yshift=-4pt]r.south west) and ([yshift=-3pt]r.south)..([xshift=-2pt,yshift=-2pt]r.south east);}. Si $x,y$ $(x\ne y)$ est un couple de points d'un polygone $\Gamma$ il existe \margin{$\cdot$}un seul co\uline{up}le d'arcs de polygone $\Gamma^{1}_{xy}$ et $\Gamma^{2}_{xy}$ d'\unsure{extremités|extrémités} $x,y$. tels que : $\mathrm{I})\quad S(\Gamma^{i}_{xy})-\{x,y\}\subset S(\Gamma),\ i=1,2$ ; $\mathrm{II})$ si $\unsure{s|z|g}$ \deleted{$\unsure{(\Gamma_{xy})|(x,y)}$}
\margin{\unsure{faut|faux}}
% slanted
\par\noindent$\unsure{\subset|c}$
% archive sheet number: 20
